\section{Size-change graphs}

\begin{Def}[Program frame]\rm
A program frame is defined to be (\Functionnames, \Callsitenames, \arity)
where
\Functionnames, \Callsitenames\ are a finite set,
and \arity : \Functionnames $\to$ N.
\end{Def}

A program frame is an abstraction of a functional program
where \Functionnames\ is the set of function names in the program,
\arity(f) is the arity of a function $f$
and each function call is the form of $f^c$
with $f$ in \Functionnames\ and $c$ in \Callsitenames\
so that each function call $f$ is 
distinguished by
attacching a different call-site $c$ to it.

For simplicity, we write SC for "size-change".

\begin{Def}[SC call graph]\rm
For a given program frame (\Functionnames,\Callsitenames, \arity),
a SC call graph is defined to a labeled graph $(V,E)$ where

vertexes $V \subseteq \Functionnames$,

edges $E \subseteq \{(f, (S, P), c, g)) \ |\ 
f,g \in \Functionnames, c \in \Callsitenames,
S, P \subseteq [1,\arity(f)] \times [1,\arity(g)],
S \cap P = \emptyset \}$,

where the graph is a directed graph with edges and
the edge $(f, (S, P), c, g))$ is from the vertex $f$ to the vertex $g$.

For simplicity, we defined $V \subseteq \Functionnames$, but
formally $V$ 
can be defined as a finite set 
with a labelling function
\[
L : \hbox{(the finite set)} \to \Functionnames.
\]
A similar formal definition can be applied to $E$.

We will often write $G_i$ for an edge in $E$.

A path is defined as a consecutive sequence of edges in $E$ as usual.
A path may be infinite.
We call an edge in a SC call graph a {\em SC graph}.
We call a path in a SC call graph a {\em SC path}.

For a given finite SC path $(G_1,\ldots,G_k)$,
a {\em trace} in the path is defined to be
a finite sequence $(n_1,\ldots,n_k)$ of natural numbers such that
$n_i (S_i \cup P_i) n_{i+1}$ for all $i \in [1,k-1]$
where $G_i = (-,(S_i,P_i),-,-)$.
For an infinite SC path,
a trace of the path is defined similarly,
where the trace is infinite.

A trace $(n_1,\ldots,n_k)$ is called progressing
if there is some $i \in [1,k-1]$ such that
$n_i P_i n_{i+1}$.
The trace is called non-progressing otherwise.
\end{Def}

\begin{Def}[Infinitely progressing trace]\rm
An initinite trace is defined to be {\em infinitely progressing}
if 
there are infinitely many $i \in N$ such that $n_i P_i n_{i+1}$
where
the trace is a trace $n_1,n_2,\ldots$ of a infinite SC path $G_1,G_2,\dots$
and
$G_i=(-,(S_i,P_i),-,-)$ for all $i \in N$.
\end{Def}


An edge $(f, (S, P), c, g)$ corresponds to
a size-change graph $G : f \to g$ in [Ikebuchi2025]
where
\[
x_i^f {\displaystyle \mathop{\to}^\downarrow} x_j^g \hbox{ if } iPj, \\
x_i^f {\displaystyle \mathop{\to}^{=\hskip -5pt\downarrow}} x_j^g \hbox{ if } iSj.
\]

The relation $S$ means a stationary relation, and
a non-progressing trace in the definition
corresponds to a non-progressing trace in [Brotherston2011].
The relation $P$ means a progressing relation, and
a progressing trace in the definition
corresponds to a progressing trace in [Brotherston2011].

\begin{Def}[SCT condition]\rm
A SC call graph is defined to satisfy SCT condition if
for any infinite SC path $G_1,G_2,\ldots$ in the SC call graph
there are some $k \in N$ and an infinitely progressing trace $n_k,n_{k+1},\ldots$
in the SC path $G_k,G_{k+1},\ldots$.
\end{Def}

\subsection{Ranking function}

We will sometimes write a sequcne $v_1,\ldots,v_k$ for $\Vec v$.
We write $(v_1,\ldots,v_k)_i$ for $v_{i+1}$ with $i \ge 0$.

\begin{Def}[State]\rm
For a program frame (\Functionnames, \Callsitenames, \arity),
a {\em state} is defined to be $(f,\Vec v)$
where $f$ is in \Functionnames\
and $\Vec v \in N^{\arity(f)}$.

For a given SC call graph of a program frame and
a call sites $c$ of the program frame and
states $(f, \Vec v)$ and $(g, \Vec w)$ of the program frame,
we define $(f, \Vec v) \to^c (g, \Vec w)$ to hold if
there is some SC graph $(f,(S,P),c,g)$ in the SC call graph such that
for all $i \in [1,\arity(f)]$ and $j \in [1,\arity(g)]$,

- $v_i \ge w_j$ if $i S j$, and

- $v_i > w_j$ if $i P j$.
\end{Def}

When we consider a program,
$(f, \Vec v) \to^c (g, \Vec w)$ is an overapproximation of
the fact that
the definition body of $f$ contains $g^c$ and
this $g^c$ is called with the form $g^c(\Vec w)$
when this $f$ is called with the form $f(\Vec v)$.

\begin{Def}[Ranking function for SC call graph]\rm
For a given SC call graph,
a function $\rho$ is defined to be a ranking function for the SC call graph if
for some $m$, we have
$\rho: \hbox{states} \to (N \cup \{\infty\})^m$
such that

if $(f, \Vec v) \to^c (g, \Vec w)$
then $\rho(f, \Vec v) >_\Lex \rho(g, \Vec w)$

where $>_\Lex$ is the lexicographic order.
\end{Def}

\begin{Th}[Ranking function synthesis (Lee2009 th 4.11, Ben2009 th 3.3, Ben2010 th 6.17)]
If a SC call graph satisfies SCT condition,
there is an algorithm to synthesize a ranking function
for the SC call graph.
\end{Th}
Note that Ben2009 th 3.3 assumed a fan-out free condition.

The metatheory of all these papers is classical.

In this paper, 
by assuming injectivity condition, which is the same as the fun-out free condition,
we will prove the same statement
in intuitionistic logic.

